\cleardoublepage
\chapter{Wprowadzenie teoretyczne}
\label{cha:RozdzialWprowadzajacy}

W~niniejszym rozdziale przedstawiono podstawowe zagadnienia teoretyczne związane z~sortowaniem 
równoległym oraz architekturą procesorów wielordzeniowych. 
Omówiono również podstawowe miary wydajności i~efektywności obliczeń równoległych, a także zagadnienie skalowalności algorytmów.

%--------------------------------------
% początek: wskazówka
%--------------------------------------
%Każdy rozdział powinien mieć zdanie wstępu przed przejściem do podrozdziału. Rysunki powinny być wyśrodkowane, numerowane i~powinny mieć odwołania do literatury lub oznaczenie że zostały opracowane samodzielnie.
%--------------------------------------
% koniec: wskazówka
%--------------------------------------

\section{Wprowadzenie do sortowania równoległego}
\label{sec:SortowanieRownolegle}

Sortowanie jest jedną z~najczęściej wykonywanych operacji w~informatyce. 
Uporządkowane dane są łatwiejsze do dalszego przetwarzania, dlatego wiele algorytmów wymaga wcześniejszego 
posortowania danych wejściowych. W~kontekście obliczeń równoległych sortowanie ma dodatkowe znaczenie ze względu 
na jego ścisły związek z~przemieszczaniem i~dystrybucją danych pomiędzy jednostkami wykonawczymi, 
co stanowi istotny element wielu algorytmów równoległych~\cite{bib:Grama}.

Celem sortowania jest przekształcenie nieuporządkowanego ciągu elementów $S = \langle a_1, a_2, \ldots, a_n \rangle$
w~ciąg uporządkowany monotonicznie rosnąco lub malejąco. Algorytmy sortowania można podzielić na porównujące
(ang. \textit{comparison-based}) oraz nieporównujące (ang. \textit{noncomparison-based}).
Algorytmy porównujące ustalają kolejność elementów na podstawie wykonywanych między nimi porównań.
Dolne ograniczenie złożoności sekwencyjnej dla tej klasy algorytmów wynosi $\Omega(n \log n)$.
Algorytmy nieporównujące wykorzystują natomiast dodatkowe właściwości danych, takie jak ich reprezentacja
lub rozkład wartości, dzięki czemu w~określonych warunkach mogą osiągać złożoność liniową~\cite{bib:Grama,bib:CLRS4}.

Równoległe sortowanie polega na rozłożeniu danych oraz operacji potrzebnych do ich uporządkowania pomiędzy wiele
jednostek wykonawczych. Typowym podejściem jest podział zbioru wejściowego na fragmenty, ich równoległe przetwarzanie,
a~następnie odpowiednia reorganizacja lub połączenie uzyskanych wyników. Sposób realizacji poszczególnych etapów
zależy od zastosowanego algorytmu oraz modelu obliczeń równoległych~\cite{bib:Grama}.

Do najważniejszych zagadnień związanych z~projektowaniem równoległych algorytmów sortowania należą sposób podziału
danych, równomierne rozłożenie pracy, komunikacja pomiędzy jednostkami wykonawczymi oraz ich synchronizacja.
Koszt tych operacji stanowi dodatkowy narzut, który nie występuje w~takiej samej postaci w~algorytmach sekwencyjnych
i~może ograniczać korzyści wynikające ze zrównoleglenia. W~szczególności nadmierna komunikacja lub nierównomierny
podział pracy mogą prowadzić do sytuacji, w~której część jednostek wykonawczych pozostaje bezczynna, podczas gdy
pozostałe nadal wykonują obliczenia~\cite{bib:Grama,bib:Jaja}.

\begin{figure}[htbp]
    \centering \includegraphics[width=0.80\textwidth]{images/chapter2/parallelSortingGeneral.png}
    \caption{Ogólny schemat podziału pracy w~sortowaniu równoległym. Źródło: opracowanie własne.}
    \label{fig:parallel-sorting-general}
\end{figure}

\section{Charakterystyka procesorów wielordzeniowych}
\label{sec:ProcesoryWielordzeniowe}

Rozwój współczesnych architektur procesorów jest związany między innymi z~ograniczeniami dalszego zwiększania wydajności 
pojedynczego strumienia instrukcji. Możliwości wykorzystania równoległości na poziomie instrukcji 
(ang. \textit{instruction-level parallelism}, ILP) są ograniczone, ponieważ pojedynczy wątek nie zawsze jest w stanie 
w~pełni wykorzystać dostępne zasoby wykonawcze procesora. Jednym ze sposobów zwiększenia stopnia wykorzystania tych zasobów 
jest wykorzystanie równoległości na poziomie wątków (ang. \textit{thread-level parallelism}, TLP)~\cite{bib:Hennessy}.

Procesory wielordzeniowe zawierają wiele jednostek obliczeniowych zdolnych do równoczesnego wykonywania różnych zadań. 
Poszczególne rdzenie mogą wykonywać niezależne strumienie instrukcji, co umożliwia jednoczesne realizowanie wielu zadań. 
Współczesne procesory mogą dodatkowo wykorzystywać mechanizmy wielowątkowości sprzętowej, takie jak 
simultaneous multithreading (SMT). Przykładem takiego rozwiązania jest technologia Intel Hyper-Threading, 
w której jeden rdzeń fizyczny udostępnia dwa konteksty wykonawcze, widoczne przez system operacyjny jako dwa procesory logiczne
~\cite{bib:IntelHT,bib:Hennessy}.

Procesory logiczne nie są jednak odpowiednikami dodatkowych niezależnych rdzeni fizycznych, ponieważ współdzielą część zasobów wykonawczych 
jednego rdzenia. Technologia SMT pozwala lepiej wykorzystać zasoby procesora w sytuacji, gdy pojedynczy wątek nie jest w stanie 
w pełni wykorzystać dostępnych jednostek wykonawczych. Uzyskiwany wzrost wydajności zależy jednak od charakterystyki obciążenia 
i~nie musi być proporcjonalny do liczby obsługiwanych wątków sprzętowych~\cite{bib:Hennessy}.

W systemach z~pamięcią współdzieloną wiele jednostek wykonawczych korzysta ze wspólnej fizycznej pamięci operacyjnej. 
Dane współdzielone mogą być wykorzystywane do komunikacji pomiędzy jednostkami poprzez operacje odczytu i zapisu w tej samej przestrzeni pamięci. 
Wykorzystanie pamięci współdzielonej wiąże się jednak z dodatkowymi wymaganiami związanymi z~organizacją hierarchii pamięci 
oraz~utrzymaniem spójnego widoku danych przechowywanych w pamięciach podręcznych poszczególnych rdzeni~\cite{bib:Hennessy,bib:Grama}.

Istotnym problemem w~architekturach z pamięcią współdzieloną jest spójność pamięci podręcznych 
(ang. \textit{cache coherence}). Te same dane mogą znajdować się jednocześnie w pamięciach podręcznych kilku rdzeni, 
dlatego zmiana wartości przez jedną jednostkę wykonawczą musi zostać odpowiednio uwzględniona przez pozostałe. 
Mechanizmy utrzymywania spójności mogą generować dodatkowy ruch pomiędzy pamięciami podręcznymi oraz wpływać na wydajność programów równoległych
~\cite{bib:Hennessy}.

Wydajność procesora wielordzeniowego zależy nie tylko od liczby dostępnych rdzeni, lecz również od organizacji systemu pamięci. 
Do ważnych parametrów należą opóźnienie dostępu do pamięci (ang. \textit{latency}) oraz~jej przepustowość (ang. \textit{bandwidth}). 
W~przypadku aplikacji intensywnie korzystających z pamięci ograniczeniem wydajności może być szybkość dostarczania danych 
do jednostek wykonawczych, a nie sama moc obliczeniowa procesora~\cite{bib:Grama}.

Zwiększenie liczby rdzeni zwiększa dostępny poziom równoległości sprzętowej, jednak rzeczywisty wzrost wydajności 
zależy również od organizacji systemu pamięci, kosztów synchronizacji oraz sposobu wykorzystania dostępnych zasobów. 
Większa liczba rdzeni nie musi więc prowadzić do proporcjonalnego skrócenia czasu wykonania~\cite{bib:Grama,bib:Hennessy}.

\section{Wydajność i efektywność}
\label{sec:WydajnoscEfektywnosc}

Podstawowymi metrykami wykorzystywanymi do oceny algorytmów równoległych są czas wykonania,
przyspieszenie (ang. \textit{speedup}) oraz efektywność (ang. \textit{efficiency}).
Umożliwiają one określenie, w~jakim stopniu zastosowanie wielu jednostek wykonawczych wpływa na czas wykonania
zadania oraz jak skutecznie wykorzystywane są dostępne zasoby obliczeniowe.

Niech $T_1(n)$ oznacza czas wykonania wariantu sekwencyjnego algorytmu sekwencyjnego dla danych o~rozmiarze $n$, 
natomiast $T_p(n)$ oznacza czas wykonania wariantu równoległego przy wykorzystaniu $p$ jednostek wykonawczych. 
Przyspieszenie definiuje się jako:

% \[
% S_p(n) = \frac{T_1(n)}{T_p(n)}.
% \]
\begin{equation}
    S_p(n) = \frac{T_1(n)}{T_p(n)}.
\end{equation}

W przypadku idealnym przyspieszenie powinno być zbliżone do liczby
wykorzystanych jednostek wykonawczych, czyli $S_p(n) \approx p$. W praktyce wartość ta może być mniejsza między innymi 
ze względu na ograniczony poziom współbieżności, synchronizację, komunikację oraz narzuty związane 
z realizacją obliczeń równoległych~\cite{bib:Jaja,bib:Grama}.

Efektywność określa stopień wykorzystania dostępnych jednostek wykonawczych i~jest definiowana
jako stosunek uzyskanego przyspieszenia do ich liczby:

% \[
% E_p(n) = \frac{S_p(n)}{p}.
% \]
\begin{equation}
    E_p(n) = \frac{S_p(n)}{p}.
\end{equation}

Po uwzględnieniu definicji przyspieszenia zależność tę można zapisać również w~postaci:

% \[
% E_p(n) = \frac{T_1(n)}{p \cdot T_p(n)}.
% \]
\begin{equation}
    E_p(n) = \frac{T_1(n)}{p \cdot T_p(n)}.
\end{equation}

Wartość efektywności bliska $1$ oznacza wysokie wykorzystanie dostępnych zasobów obliczeniowych.
W rzeczywistych systemach efektywność jest zazwyczaj niższa od wartości idealnej ze względu na czas
poświęcany między innymi na komunikację, synchronizację oraz oczekiwanie jednostek wykonawczych
~\cite{bib:Jaja,bib:Grama}.

Istotnym ograniczeniem możliwego przyspieszenia jest udział części programu, która nie może zostać
wykonana równolegle. Zależność tę opisuje prawo Amdahla. W jego podstawowym modelu czas wykonania
programu dzielony jest na część sekwencyjną oraz część możliwą do wykonania równolegle.
Czas części sekwencyjnej pozostaje niezależny od liczby wykorzystywanych jednostek wykonawczych,
natomiast czas idealnie skalowalnej części równoległej zmniejsza się proporcjonalnie do ich liczby~\cite{bib:KarbowskiAmdahl}.

Jeżeli przez $f$ oznaczony zostanie udział części sekwencyjnej programu, a przez $p$ liczbę
jednostek wykonawczych, przyspieszenie wynikające z prawa~Amdahla można zapisać jako:

% \[
% S_p = \frac{1}{f + \frac{1-f}{p}}.
% \]
\begin{equation}
    S_p = \frac{1}{f + \frac{1-f}{p}}.
\end{equation}

Zależność ta wskazuje, że nawet przy zwiększaniu liczby jednostek wykonawczych część sekwencyjna
programu ogranicza maksymalne możliwe przyspieszenie. Klasyczny model prawa Amdahla zakłada przy tym
idealne skalowanie części równoległej i~pomija dodatkowe koszty występujące w rzeczywistych systemach.
Do takiego narzutu zalicza się komunikację, synchronizację oraz tworzenie procesów lub wątków~\cite{bib:KarbowskiAmdahl}.

\section{Skalowalność algorytmów równoległych}
\label{sec:SkalowalnoscAlgorytmow}

Skalowalność systemu równoległego określa, jak efektywnie wykorzystuje on rosnącą liczbę jednostek wykonawczych 
i~jaki wzrost przyspieszenia uzyskuje wraz ze zwiększaniem ich liczby. 
Dobrze skalowalny system powinien umożliwiać zwiększanie liczby jednostek wykonawczych bez gwałtownego spadku efektywności~\cite{bib:Grama}.

Przy stałym rozmiarze danych wejściowych zwiększanie liczby jednostek wykonawczych zazwyczaj powoduje stopniowe obniżanie efektywności obliczeń. 
Powodem tego jest wzrost całkowitego narzutu obliczeń równoległych, na który składają się takie czynniki jak komunikacja pomiędzy jednostkami, 
czas ich bezczynności oraz dodatkowe obliczenia wykonywane w wariancie równoległym. 
Wzrost liczby jednostek wykonawczych może przyczyniać się do zwiększenia tych kosztów, 
co skutkuje ograniczeniem uzyskiwanego przyspieszenia~\cite{bib:Grama}.

Skalowalność algorytmu zależy również od rozmiaru przetwarzanych danych.
Dla wielu systemów równoległych zwiększenie rozmiaru danych wejściowych przy stałej liczbie jednostek wykonawczych powoduje wzrost efektywności, 
ponieważ większy udział w całkowitym czasie wykonania mają właściwe obliczenia. 
System równoległy uznaje się za skalowalny, jeżeli przy jednoczesnym zwiększaniu liczby jednostek wykonawczych 
oraz rozmiaru danych możliwe jest utrzymanie efektywności na zbliżonym poziomie~\cite{bib:Grama}.

Jedną z miar wykorzystywanych do oceny skalowalności jest funkcja \textit{isoefficiency}, 
przedstawiająca zależność między wzrostem liczby jednostek wykonawczych $p$ a rozmiarem danych $W$, 
niezbędnym do zachowania stałego poziomu efektywności.
Im wolniej musi rosnąć rozmiar zadania obliczeniowego wraz ze wzrostem $p$, tym lepszą skalowalnością charakteryzuje się dany system równoległy. 
Niższe tempo wzrostu funkcji \textit{isoefficiency} oznacza więc możliwość efektywnego wykorzystania 
większej liczby jednostek wykonawczych przy stosunkowo niewielkiej zmianie rozmiaru zadania obliczeniowego~\cite{bib:Grama}.

Dolną granicą funkcji \textit{isoefficiency} jest wzrost liniowy względem liczby jednostek wykonawczych. 
Idealnie skalowalny system charakteryzuje się zatem funkcją \textit{isoefficiency} rzędu $\Theta(p)$, co oznacza, 
że dla zachowania stałej efektywności wystarczający jest liniowy wzrost rozmiaru danych wejściowych 
wraz ze wzrostem liczby jednostek wykonawczych~\cite{bib:Grama}.

Analiza skalowalności umożliwia ocenę, w~jaki sposób system równoległy zachowuje się przy 
zwiększaniu liczby jednostek wykonawczych oraz rozmiaru przetwarzanych danych.